\section{Two cut-square theorems}\label{sec:metric-cut-square}

The phrase \emph{cut square} appears in two mathematically distinct identities. Their separation is essential.

\subsection{Operator cut square}

Under the cut-odd generator hypothesis, \cref{thm:t41bilateral} gives
\[
 \boxed{(U_t+U_{-t})^2-(U_t-U_{-t})^2=4I.}
\]
This is an identity in the operator algebra generated by $G$.

\subsection{Positive-metric cut square}

Let $\mathcal R\cong\C^m$ be a finite response fibre and let $H>0$ be a declared metric. For $y,c\in\mathcal R$, define
\[
 \langle u,v\rangle_H=u^*Hv,
 \qquad
 \|u\|_H^2=u^*Hu,
\]
\[
 L_c(y)=\langle c,y\rangle_H,
 \qquad
 \beta_H(c)=\|c\|_H^2.
\]
Put $C=H^{1/2}c$ and $Y=H^{1/2}y$.

\begin{theorem}[Positive-metric cut-square identity]\label{thm:metriccutsquare}
For every response event $y$ and decoder $c$,
\[
 \boxed{
 \beta_H(c)\|y\|_H^2-|L_c(y)|^2
 =\sum_{i<j}|C_iY_j-C_jY_i|^2
 =\|C\wedge Y\|^2.}
\]
If $\beta_H(c)\le1$, then
\[
 \boxed{
 \|y\|_H^2-|L_c(y)|^2
 =(1-\beta_H(c))\|y\|_H^2+\|C\wedge Y\|^2\ge0.}
\]
Equivalently,
\[
 H-Hcc^*H\ge0.
\]
\end{theorem}

\begin{proof}
In whitened coordinates, the complex Lagrange identity gives
\[
 \|C\|^2\|Y\|^2-|\langle C,Y\rangle|^2
 =\sum_{i<j}|C_iY_j-C_jY_i|^2.
\]
Substitution proves the first formula. Add $(1-\beta_H(c))\|y\|_H^2$ to both sides to obtain the contractive decomposition. The operator inequality is the same statement at the quadratic-form level.
\end{proof}

\begin{corollary}[Equality classification]
If $\beta_H(c)<1$, equality in the contractive decomposition occurs only for $y=0$. If $\beta_H(c)=1$, equality occurs exactly when $Y$ is a scalar multiple of $C$.
\end{corollary}

\subsection{Four-channel tetrad}

For four normalized response channels $p,v,s,t$, the transverse residue is the sum of six squares:
\[
\begin{aligned}
\|C\wedge Y\|^2={}&|C_pY_v-C_vY_p|^2+|C_pY_s-C_sY_p|^2
+|C_pY_t-C_tY_p|^2\\
&+|C_vY_s-C_sY_v|^2+|C_vY_t-C_tY_v|^2
+|C_sY_t-C_tY_s|^2.
\end{aligned}
\]
One scalar closure relation can hold while cross-channel minors remain nonzero. The six-minor condition is therefore a finer coherence test than a single ratio identity.

\subsection{Relation, without conflation}

The operator identity compares forward and backward functions of one generator. The metric identity compares a response event with one decoder direction. The first uses algebraic inverse flow; the second uses positivity and an inner product. An adapter may place both identities in one model, but neither proves the other.
